% HCLT 2026 양식 (Chanjun Park, Cheoneum Park)
% Overleaf 사용 시 Menu -> TeX Live version -> 2021 로 설정
%
% 주: 템플릿 원본의 \usepackage[utf8]{inputenc} 는 제거하였다.
%     LaTeX2e 는 2018-04-01 부터 UTF-8 이 기본 입력 인코딩이므로 불필요하며,
%     hclt.sty 가 kotex 을 엔진 감지 없이 불러오는 탓에 XeLaTeX/LuaLaTeX 에서
%     inputenc 가 충돌한다. 제거하면 세 엔진 모두 통과한다.

\documentclass[twocolumn, 10pt]{article}

\usepackage{hclt}
\usepackage[pdfencoding=auto]{hyperref}
\usepackage{url}
\usepackage{amsmath}
\usepackage{latexsym}
\usepackage{amssymb}
\usepackage{graphicx}

\DeclareMathOperator*{\argmax}{arg\,max}
\DeclareMathOperator*{\argmin}{arg\,min}

\title[korean]{한국어 학습 데이터 없이 영어 탈옥 탐지기를 한국어로 전이하기: \\ 중간층 표현 정렬을 이용한 언어 간 성능 격차 해소}
\title[english]{Transferring English Jailbreak Detectors to Korean without Korean Training Data: \\ Closing the Cross-Lingual Gap via Middle-Layer Representation Alignment}

\author[korean]{
홍길동$^{\circ}$, 김길수 \\
한국대학 한국학부\\
\texttt{\{hong,kim\}@mail.ac.kr}\\
}

\author[english]{
Gil-Dong Hong$^{\circ}$, Gil-Su Kim \\
Hanguk University, Hanguk Research Institute
}

% \finalcopy % Uncomment for camera-ready version

\begin{document}

\twocolumn[
\maketitle
\begin{abstract}
한국어 탈옥(jailbreak) 프롬프트 데이터셋과 이를 탐지하는 경량 분류기는 공개된 것이 사실상 없다.
이에 실무에서는 영어로 학습된 공개 가드레일을 한국어에 그대로 적용하지만, 그 성능은 지금까지 보고된 바 없다.
본 논문은 영어 탈옥 벤치마크 세 종을 한국어로 번역해 영·한 병렬 평가셋을 구축하고, 공개 가드레일의 한국어 성능을 최초로 측정한다.
측정 결과 Meta의 Prompt Guard 2는 오탐률 1\% 기준 재현율이 영어 0.4756에서 한국어 0.2711로 20.5\%p 하락한다.
층별 번역 검색 분석으로 원인이 상위층에서 언어 간 표현 정렬이 붕괴하는 데 있음을 보이고(8층 93.00\%, 12층 3.33\%), 이 붕괴가 미세조정이 아니라 사전학습 구조에서 기인함을 확인한다.
이에 중간층 대조 정렬 손실을 과제 학습에 결합한 방법을 제안한다.
제안 방법은 한국어 라벨을 전혀 사용하지 않고 영·한 격차를 0.2\%p로 줄이며, 한국어 재현율에서 교사 모델을 6.1\%p 앞선다($p<10^{-13}$).
\end{abstract}

\begin{keyword}
탈옥 탐지, 가드레일, 교차언어 전이, 표현 정렬, 지식 증류
\end{keyword}
]

\section{서론}

대규모 언어모델의 배포가 확산되면서 안전장치를 우회하려는 탈옥 프롬프트가 실질적 위협이 되었다.
이를 막기 위해 Meta의 Prompt Guard, ProtectAI 등 경량 분류기 형태의 가드레일이 공개되었으나, 이들 모델은 모두 영어를 중심으로 학습·평가되었다.
Prompt Guard v1과 v2는 모델 카드에 평가 언어를 영어, 프랑스어, 독일어, 힌디어, 이탈리아어, 포르투갈어, 스페인어, 태국어 여덟 가지로 명시하고 있으며 한국어는 포함되어 있지 않다.

한국어 상황은 더 열악하다.
공개된 한국어 안전성 자원은 KoSBi, SQuARe와 같은 사회적 편향·민감 주제 데이터셋이나 MultiJail~\cite{deng:iclr2024}, PolyGuard~\cite{kumar:colm2025}처럼 유해 요청을 다루는 자료가 대부분이며, 탈옥 프롬프트 자체를 대상으로 한 학회 검증 데이터셋은 확인되지 않는다.
국내에서는 영어-한국어 탈옥 프롬프트 데이터셋을 구축한 연구가 보고되었으나~\cite{park:kiisc2025}, 기술이전 계약으로 모델과 데이터가 공개되지 않아 후속 연구가 이를 활용할 수 없다.

이러한 공백 때문에 실무에서는 영어 가드레일을 한국어 입력에 그대로 적용하는 것이 사실상 유일한 선택지가 된다.
그러나 그 방식이 실제로 얼마나 동작하는지는 측정된 적이 없다.
본 논문은 이 질문에서 출발하여 다음 세 가지를 기여한다.

\begin{itemize}
\item 영어 탈옥 벤치마크 세 종을 번역해 영·한 병렬 평가셋을 구축하고, 공개 가드레일의 한국어 탈옥 탐지 성능을 최초로 보고한다.
\item 층별 번역 검색 분석으로 언어 간 성능 격차의 원인이 상위층 표현 정렬 붕괴임을 규명하고, 이 붕괴가 사전학습 단계에서 이미 존재함을 보인다.
\item 중간층 대조 정렬 손실을 제안하여, 한국어 라벨 없이 영·한 격차를 0.2\%p로 줄인다.
\end{itemize}

\section{관련 연구}

\subsection{교차언어 전이}
XTREME~\cite{hu:icml2020}은 다국어 전이 평가를 zero-shot, translate-train, translate-test 세 설정으로 정리하였다.
Lauscher 등~\cite{lauscher:emnlp2020}은 zero-shot 전이가 저자원 언어에서 한계를 보이며 소량의 목표어 사례가 이를 상당 부분 보완함을 보였다.
K 등~\cite{k:iclr2020}은 통제 실험을 통해 어휘 중첩이 교차언어 능력에 미치는 영향은 미미하고 신경망의 깊이가 핵심 요인임을 밝혔다.

\subsection{중간층 표현 정렬}
Liu와 Niehues~\cite{liu:acl2025}는 1{,}000개 이상의 언어쌍에 대한 분석에서 중간층이 교차언어 정렬 잠재력이 가장 크다는 점을 발견하고, 중간층에 대조 정렬 목적함수를 결합하는 방법을 제안하였다.
해당 연구는 생성형 대규모 언어모델을 대상으로 슬롯 채우기, 기계번역, 구조화 텍스트 생성 과제에서 효과를 검증하였다.
본 논문은 이 아이디어를 인코더 기반 분류기와 안전 가드레일 과제에 적용하고, 한국어를 단일 목표 언어로 하여 층 선택의 영향을 분석한다.

\subsection{다국어 안전 가드레일}
PolyGuard~\cite{kumar:colm2025}는 17개 언어를 지원하는 안전 조정 모델과 191만 건 규모의 학습 자료를 공개하였고, 한국어를 포함한다.
CSRT~\cite{yoo:acl2025}는 코드 스위칭 질의로 다국어 안전성을 평가한다.
그러나 두 연구 모두 유해 콘텐츠 판정을 과제로 삼으며, 한국어 성능을 개별적으로 보고하지 않는다.
PIGuard~\cite{li:acl2025}는 트리거 단어에 의한 과탐 문제를 다루며 NotInject 평가셋을 제안하였는데, 본 논문은 이를 한국어로 번역해 오탐 측정에 활용한다.

\section{한국어 탈옥 평가셋 구축}

\subsection{원본 벤치마크}
과제는 Prompt Guard 2 모델 카드의 정의를 따라 \emph{이전 지시나 안전장치를 덮어쓰려는 시도}의 탐지로 한정한다.
이 정의에 부합하는 영어 벤치마크 세 종을 선정하였다.

\begin{itemize}
\item \textbf{JailbreaksOverTime}: 2023년 2월부터 12월까지 수집된 실사용 프롬프트 23{,}558건. 시간 분할을 적용해 2023년 12월분을 평가에 사용한다.
\item \textbf{in-the-wild jailbreak prompts}~\cite{shen:ccs2024}: 온라인 커뮤니티에서 실제로 유포된 탈옥 템플릿.
\item \textbf{WildJailbreak}~\cite{jiang:neurips2024}: WildTeaming 프레임워크가 합성한 적대적 탈옥 프롬프트.
\end{itemize}

정상 표본은 JailbreaksOverTime의 실사용 발화와 NotInject~\cite{li:acl2025} 두 출처에서 구성하였다.
NotInject는 ``ignore'', ``roleplay''와 같이 탈옥 프롬프트가 즐겨 쓰는 단어를 포함하지만 무해한 문장으로, 어휘 기반 지름길을 차단하는 역할을 한다.

\subsection{번역 및 품질 관리}
번역기는 TowerInstruct-7B-v0.2~\cite{alves:emnlp2024}를 사용하였다.
PolyGuard~\cite{kumar:colm2025}가 PolyGuardMix 구축에 사용한 모델과 동일하며, 한국어를 지원 언어에 포함한다.
긴 프롬프트는 문장 경계로 900자 이내 조각으로 분할해 번역한 뒤 이어 붙였고, 반복 붕괴를 억제하기 위해 반복 벌점 1.08과 12-gram 반복 금지를 적용하였다.

번역 결과에는 네 가지 자동 검사를 적용하였다.
첫째, 한글·라틴 문자 외의 문자가 혼입된 경우를 탐지하였다.
TowerInstruct가 10개 언어 모델인 탓에 키릴 문자나 한자가 섞이는 사례가 관측되었으며(291건 중 15건), 청크 크기를 변경한 재번역과 문자 수준 교정으로 모두 제거하였다.
둘째, 원문 대비 길이 비가 0.15에서 2.5 범위를 벗어난 항목을 제외하였다.
셋째, 한글이 전혀 없는 항목을 제외하였다.
넷째, 다국어 문장 임베딩 모델로 원문과 번역문의 교차언어 유사도를 측정하였다.
짝이 맞는 쌍의 유사도 중앙값은 0.886이었고 무작위로 섞은 쌍은 0.539로, 변별력은 +0.32였다.

\subsection{근접 중복 측정}
탈옥 프롬프트는 동일한 템플릿이 여러 달에 걸쳐 재유포되는 특성이 있다.
정확 일치 기준으로 중복을 제거하더라도 변형본이 남을 수 있으므로, 학습셋과 평가셋 사이의 문자 5-gram 자카드 유사도를 측정하였다.
JailbreaksOverTime 2023년 12월 시험셋의 공격 286건 중 233건(81.5\%)이 학습셋에 유사도 0.7 이상인 문장을 가지고 있었다.
따라서 시간 분할만으로는 암기 효과를 완전히 배제할 수 없으며, 본 논문은 이 지표를 함께 보고한다.

\subsection{최종 구성}
표 \ref{tab:eval}\은 최종 평가셋의 구성이다.
영어판과 한국어판은 동일한 항목을 동일한 순서로 담고 있어 언어만을 변수로 하는 비교가 가능하다.
임계값 선택용 검증셋은 정상 300건으로 별도 구성하였다.

\begin{table}[t]
\caption{한국어 탈옥 평가셋 구성}
\label{tab:eval}
\begin{center}
\begin{tabular}{l|l|r}
\hline
구분 & 출처 & 건수 \\
\hline
공격 & WildJailbreak & 1,168 \\
     & in-the-wild & 394 \\
     & JailbreaksOverTime & 286 \\
\hline
정상 & JailbreaksOverTime & 467 \\
     & NotInject & 170 \\
\hline
\multicolumn{2}{c|}{합계} & 2,485 \\
\hline
\end{tabular}
\end{center}
\end{table}

\section{언어 간 성능 격차의 원인 분석}

\subsection{층별 번역 검색}
언어 간 표현이 어느 층에서 만나는지 확인하기 위해 번역 검색을 진행하였다.
같은 의미의 영·한 문장 1{,}200쌍에 대해 각 층의 은닉 상태를 어텐션 마스크로 가중 평균하여 문장 벡터를 얻은 뒤, 한국어 문장 벡터와 가장 가까운 영어 문장 벡터가 실제 번역 짝인지를 판정하였다.
무작위 선택의 기댓값은 0.08\%이다.

표 \ref{tab:layer}\는 그 결과이다.
사전학습만 거친 mDeBERTa-v3-base~\cite{he:iclr2023}는 8층에서 93.00\%를 기록하나 11층에서 2.17\%로 급락한다.
영어 전용 DeBERTa-v3는 최고값이 40.33\%에 그쳐, 다국어 사전학습이 중간층 정렬을 만든다는 점을 보여준다.

\begin{table}[t]
\caption{층별 번역 검색 정확도(\%). 1,200개 후보 중 정답 선택 비율}
\label{tab:layer}
\begin{center}
\begin{tabular}{c|c|c|c}
\hline
층 & mDeBERTa & 영어 DeBERTa & Prompt Guard 2 \\
\hline
0  & 8.50  & 6.50  & 8.42 \\
4  & 30.83 & 10.83 & 26.00 \\
6  & 87.75 & 31.17 & 79.42 \\
8  & \textbf{93.00} & 40.33 & \textbf{81.83} \\
10 & 81.75 & 14.58 & 53.50 \\
11 & 2.17  & 5.25  & 3.25 \\
12 & 3.33  & 0.92  & 6.50 \\
\hline
\end{tabular}
\end{center}
\end{table}

\subsection{붕괴의 원인}
상위층 붕괴가 영어 과제 미세조정에서 비롯되었을 가능성을 검토하였다.
미세조정 전 원본 모델에서 이미 11층이 2.17\%였고, 영어 데이터로 미세조정한 뒤에는 오히려 3.33\%에서 4.00\%로 소폭 상승하였다.
따라서 붕괴는 미세조정의 결과가 아니다.

다국어 표현이 언어 성분과 의미 성분으로 분해된다는 관찰~\cite{libovicky:csl2020}과 부합하며, 원인은 사전학습 목적함수에 있는 것으로 해석된다.
DeBERTa-v3는 치환된 토큰을 판별하도록 학습되며, 그 판정은 언어별 어휘 단위로 이루어져야 한다.
출력에 가까운 층일수록 언어별 표면형으로 되돌아가는 것이 목적함수를 만족시키는 방향이다.

이 붕괴는 대규모 다국어 학습으로도 해소되지 않는다.
여덟 개 언어로 학습된 Prompt Guard 2 역시 11층에서 3.25\%를 기록한다.

\section{제안 방법}

\subsection{구성}
교사 모델로 Prompt Guard 2 86M을 사용한다.
백본이 학생 모델과 동일한 mDeBERTa-v3-base이며, 학습 언어 여덟 개에 한국어가 포함되지 않는다.
교사는 갱신하지 않고 영어 문장에 대한 확률만 산출한다.

학생 모델은 과제를 학습한 적이 없는 mDeBERTa-v3-base이다.
학습 자료는 JailbreaksOverTime 2023년 2월~11월 구간의 영어 프롬프트 18{,}585건이며, 원 라벨은 사용하지 않는다.
정렬용으로는 그 일부를 번역한 영·한 병렬 1{,}958쌍을 사용한다.

\subsection{목적함수}
전체 손실은 과제 손실과 정렬 손실의 합이다.
\begin{eqnarray}
L = L_{\text{task}} + \lambda L_{\text{align}}. \label{eqn:total}
\end{eqnarray}

과제 손실은 교사의 확률분포를 학생이 재현하도록 하는 교차 엔트로피이다.
\begin{eqnarray}
L_{\text{task}} = -\sum_{k} q_k(x) \log p_k(x), \nonumber
\end{eqnarray}
여기서 $q$는 교사가 영어 문장 $x$에 부여한 확률, $p$는 학생의 확률이다.

정렬 손실은 영어 문장 $s$와 그 번역 $t$의 $i$층 표현을 가깝게 하고 짝이 아닌 문장과는 멀게 하는 대조 손실이다~\cite{liu:acl2025}.
\begin{eqnarray}
L_{\text{align}} = -\log \frac{\exp(\mathrm{sim}(h^i_s, h^i_t)/\tau)}{\sum_{v \in B} \exp(\mathrm{sim}(h^i_s, h^i_v)/\tau)}, \label{eqn:align}
\end{eqnarray}
$h^i$는 $i$층 은닉 상태의 마스크 가중 평균, $\mathrm{sim}$은 코사인 유사도, $B$는 미니배치, $\tau$는 온도이다.

수식 \eqref{eqn:align}\은 과제 라벨을 요구하지 않는다.
필요한 정보는 두 문장이 번역 관계라는 사실뿐이며, 한국어에는 어떠한 안전성 라벨도 부여되지 않는다.

\subsection{학습 설정}
3 에폭, 학습률 $2\times10^{-5}$, 배치 크기 16, $\lambda=1.0$, $\tau=0.05$로 학습하였다.
정렬 층 $i$는 4, 6, 8, 10으로 바꾸어 가며 비교하였고, 정렬 손실을 제거한 조건($\lambda=0$)을 기준선으로 삼았다.

\section{실험}

\subsection{평가 지표}
탈옥 탐지는 오탐 비용이 크므로 오탐률을 고정한 재현율로 평가한다.
임계값은 검증셋에서만 선택해 시험셋에 그대로 적용하며, 오탐률 1\%와 0.1\% 두 지점을 보고한다.

\subsection{주 결과}
표 \ref{tab:main}\은 주 결과이다.
교사 Prompt Guard 2는 영어 0.4756에서 한국어 0.2711로 20.5\%p 하락한다.
정렬 손실 없이 증류만 수행한 기준선은 9.2\%p 하락한다.
8층 정렬을 적용하면 격차가 0.2\%p로 줄어들며, 한국어 재현율은 교사를 6.1\%p 앞선다.

층 선택은 결과를 크게 좌우한다.
정렬 정확도가 낮은 4층에 손실을 걸면 기준선보다 오히려 성능이 떨어진다(0.2008 대 0.2440).
표 \ref{tab:layer}에서 정렬 정확도가 가장 높았던 8층이 전이 성능에서도 최적이었다.

\begin{table}[t]
\caption{한국어 탈옥 탐지 결과. 오탐률 1\% 고정 재현율}
\label{tab:main}
\begin{center}
\begin{tabular}{l|c|c|c}
\hline
모델 & 영어 & 한국어 & 격차 \\
\hline
Prompt Guard 2 & \textbf{0.4756} & 0.2711 & $-$0.2045 \\
\hline
정렬 없음 & 0.3355 & 0.2440 & $-$0.0915 \\
4층 정렬 & 0.2808 & 0.2008 & $-$0.0801 \\
6층 정렬 & 0.3550 & 0.3025 & $-$0.0525 \\
8층 정렬 & 0.3339 & \textbf{0.3323} & \textbf{$-$0.0016} \\
10층 정렬 & 0.3485 & 0.3252 & $-$0.0233 \\
\hline
\end{tabular}
\end{center}
\end{table}

\subsection{통계적 유의성}
한국어 공격 1{,}848건에 대해 95\% 윌슨 신뢰구간을 계산하였다.
8층 정렬은 0.3323 [0.3111, 0.3541], 교사는 0.2711 [0.2513, 0.2918], 기준선은 0.2440 [0.2250, 0.2641]로 구간이 겹치지 않는다.

동일 문장에 대한 짝지은 비교(McNemar 정확검정)에서 8층 정렬은 기준선이 놓친 181건을 추가로 탐지한 반면 그 반대는 18건에 그쳤다($p=4.7\times10^{-35}$).
교사와의 비교에서도 8층 정렬만 맞힌 사례가 175건, 교사만 맞힌 사례가 62건이었다($p=1.2\times10^{-13}$).

\subsection{공격 유형별 분석}
표 \ref{tab:src}\은 공격 출처별 재현율이다.
JailbreaksOverTime 계열은 학습 분포와 가까워 0.8986으로 높다.
반면 WildJailbreak의 합성 적대적 프롬프트는 0.1421로 가장 어렵다.
다만 이 구간에서도 정렬 손실은 기준선 대비 약 4배의 개선을 보인다.

\begin{table}[t]
\caption{공격 출처별 한국어 재현율(오탐률 1\% 고정) 및 NotInject 오탐률}
\label{tab:src}
\begin{center}
\begin{tabular}{l|c|c|c|c}
\hline
모델 & JOT & ITW & WJB & 오탐 \\
\hline
정렬 없음 & 0.8182 & 0.4416 & 0.0368 & 0.0\% \\
8층 정렬 & \textbf{0.8986} & \textbf{0.4848} & \textbf{0.1421} & 0.6\% \\
Prompt Guard 2 & 0.7972 & 0.4670 & 0.0762 & 0.6\% \\
\hline
\end{tabular}
\end{center}
\end{table}

\subsection{과탐 검증}
NotInject~\cite{li:acl2025}의 한국어 번역본에 대한 오탐률은 8층 정렬 0.6\%, 교사 0.6\%로 동일하다.
탈옥 프롬프트가 자주 사용하는 단어를 포함한 무해 문장에서 추가적인 과탐 없이 재현율 개선을 얻었음을 의미한다.

\subsection{영어 전용 백본의 한계}
백본의 다국어 여부가 전이 가능성을 결정한다는 점을 확인하기 위해 영어 전용 백본 모델을 함께 측정하였다.
PIGuard~\cite{li:acl2025}는 영어 0.8706에서 한국어 0.1503으로, ProtectAI v2는 0.4930에서 0.0000으로 하락하였다.
어휘 분석에서도 mDeBERTa는 2음절 이상 한글 토큰을 1{,}846개 보유한 반면 영어 DeBERTa-v3는 80개에 그쳤다.

\subsection{어휘 지름길 검증}
평가셋이 표면적 단서만으로 풀리는지 확인하기 위해 문자 3-gram TF-IDF 기반 선형 분류기를 대조군으로 두었다.
대조군은 학생 모델과 동일하게 영어 학습셋으로만 학습하고 한국어를 보지 않았다.
그 결과 대조군의 한국어 재현율은 0.4266에 그쳐 학생 모델(0.8287)에 크게 못 미쳤다.
또한 번역문에 남은 라틴 문자를 모두 제거하면 대조군은 0.0455로 붕괴한 반면 학생 모델은 0.8478로 변화가 거의 없었다.
어휘 규칙은 언어를 넘지 못하는 반면 표현 수준의 지식은 전이됨을 보여준다.
길이만을 점수로 사용하는 대조군은 오탐률 1\% 지점에서 재현율 0.0000을 기록하였다.

\section{한계}

첫째, 평가셋의 공격 프롬프트는 모두 영어 원문의 기계번역이며 원어민이 작성한 한국어 탈옥 프롬프트가 아니다.
자모 분해와 같이 한국어에만 존재하는 공격 유형은 포함되지 않는다.

둘째, 학습셋과 평가셋 사이에 근접 중복이 존재한다.
JailbreaksOverTime 구간에서는 공격의 81.5\%가 학습셋에 유사 문장을 가진다.
근접 중복을 제외한 부분집합에서는 모든 모델의 성능이 하락하나 모델 간 순위는 유지되었다.

셋째, 단일 시드 결과이며 반복 실험을 수행하지 않았다.

넷째, 제안 방법은 교사가 아는 지식을 한국어로 옮기는 데 효과가 있으며, 교사가 다루지 않는 과제에는 도움이 되지 않는다.
유해 요청 판정을 다루는 MultiJail~\cite{deng:iclr2024} 한국어 부분집합에서는 교사와 학생 모두 0.02 이하를 기록하였다.

\section{결론}

본 논문은 한국어 탈옥 탐지 자원이 부재한 상황에서 영어 가드레일의 한국어 성능을 최초로 측정하고, 20.5\%p의 언어 간 격차가 존재함을 보였다.
층별 번역 검색 분석으로 그 원인이 상위층 표현 정렬 붕괴이며 사전학습 구조에서 기인함을 규명하였다.
중간층 대조 정렬 손실을 결합한 제안 방법은 한국어 라벨을 전혀 사용하지 않고 격차를 0.2\%p로 줄였고, 한국어 재현율에서 교사 모델을 통계적으로 유의하게 앞섰다.
정렬 층의 선택이 결과를 좌우하며, 번역 검색 정확도가 가장 높은 층이 전이 성능에서도 최적이라는 점을 확인하였다.

향후 연구로는 원어민이 작성한 한국어 탈옥 프롬프트 수집, 여러 층에 대한 동시 정렬, 그리고 유해 콘텐츠 가드레일과의 결합을 고려한다.

\bibliographystyle{IEEEtran}
\bibliography{refer}

\end{document}
